% Preamble requirements:
% \usepackage{amsmath,amssymb,booktabs,array,enumitem,graphicx,tikz}
% \usetikzlibrary{arrows.meta,positioning,shapes.geometric}
%
% Bibliography files used by this section:
%   the proposal references.bib
%   doc/related_work_additions.bib
%   doc/proposed_methodology_additions.bib

\providecommand{\methodfield}[1]{\texttt{\detokenize{#1}}}

\section{Proposed Methodology}
\label{sec:proposed_methodology}

This section presents the proposed hierarchical intrusion-detection method in
implementation-level detail. The detector models two complementary structures
in network traffic. Stage~1 represents the ordered packets and statistical
properties of an individual bidirectional flow, whereas Stage~2 models the
relationship between a target flow and eligible historical flows. The two
stages are connected by a fixed-dimensional Stage~1 embedding and are trained
sequentially. This design makes it possible to quantify the contribution of
intra-flow timing, flow statistics, and inter-flow context independently.

The method builds on the Transformer architecture \cite{Vaswani2017}, but its
contribution is not the use of self-attention alone. It combines (i) an explicit
cumulative-time coordinate for irregular packet arrivals, (ii) a separately
encoded flow-statistics branch with bounded feature-wise token modulation, and
(iii) a causal target-query context encoder whose current flow attends only to
historical flows. Time-aware packet modelling is motivated by prior NIDS work
showing that packet order and elapsed time carry different information
\cite{Han2023GTID}. The proposed hierarchy additionally separates within-flow
evidence from across-flow behavior, following the broader observation that
independently classified flow records can miss attacks whose evidence is spread
over several connections
\cite{LiuPatras2022NetSentry,Manocchio2024FlowTransformer,
MartinReizabal2026FlowWindow}.

\subsection{Architectural Overview}
\label{subsec:method_architecture}

Let flow $i$ contain a selected and padded packet-feature matrix
$\mathbf{X}_i\in\mathbb{R}^{L\times d_p}$, a packet inter-arrival tensor
$\mathbf{u}_i\in\mathbb{R}^{L}$, a validity mask
$\mathbf{m}_i\in\{0,1\}^{L}$, a statistical flow vector
$\mathbf{s}_i\in\mathbb{R}^{d_f}$, and a binary target
$y_i\in\{0,1\}$. The principal configuration uses $L=128$. Stage~1 computes

\begin{equation}
    \left(\boldsymbol{\ell}^{(1)}_i,\mathbf{z}^{(1)}_i\right)
    =F_{\theta_1}
    \left(
        \mathbf{X}_i,\mathbf{u}_i,\mathbf{m}_i,\mathbf{s}_i
    \right),
    \label{eq:method_stage1_mapping}
\end{equation}

where $\boldsymbol{\ell}^{(1)}_i\in\mathbb{R}^{2}$ are Stage~1 logits and
$\mathbf{z}^{(1)}_i\in\mathbb{R}^{d}$ is the fused representation exported to
Stage~2. In the principal model, $d=128$.

After all Stage~1 embeddings have been exported, the flows are stably sorted by
start time and flow identifier. For a context relation $r$ and target flow $i$,
Stage~2 constructs a bounded causal sequence

\begin{equation}
    \mathbf{Z}^{(r)}_i=
    \left[
        \mathbf{z}^{(1)}_{j_1},\ldots,
        \mathbf{z}^{(1)}_{j_{K_i}},
        \mathbf{z}^{(1)}_i
    \right],
    \qquad j_k<i,quad K_i\leq W-1,
    \label{eq:method_stage2_input}
\end{equation}

where the current flow is always the last valid token and the principal context
window is $W=128$, including the target. Stage~2 predicts

\begin{equation}
    \boldsymbol{\ell}^{(2,r)}_i
    =G_{\theta_2}
    \left(\mathbf{Z}^{(r)}_i,\boldsymbol{\mu}_i\right),
    \label{eq:method_stage2_mapping}
\end{equation}

where $\boldsymbol{\mu}_i$ masks left padding. The final malicious-flow
probability is the class-1 softmax output of
$\boldsymbol{\ell}^{(2,r)}_i$.

Figure~\ref{fig:method_architecture} shows the information path. Flow labels
and Suricata alert fields are targets only and do not enter either encoder.

\begin{figure}[t]
    \centering
    \begin{tikzpicture}[
        node distance=4.5mm,
        block/.style={draw, rectangle, rounded corners=2pt, align=center,
                      text width=11.4cm, minimum height=7.5mm, fill=black!3,
                      font=\small},
        side/.style={draw, rectangle, rounded corners=2pt, align=center,
                     text width=3.2cm, minimum height=8mm, fill=black!3,
                     font=\small},
        arrow/.style={-{Latex[length=2mm]},thick}
    ]
        \node[block] (packet) {Selected packet features, inter-arrival times, and packet mask};
        \node[block, below=of packet] (projection) {Record projection and cumulative time-aware sinusoidal encoding};
        \node[block, below=of projection] (film) {Bounded flow-conditioned token modulation};
        \node[side, right=7mm of film] (flow) {Encoded flow\\statistics};
        \node[block, below=of film] (encoder) {Masked intra-flow Transformer encoder};
        \node[block, below=of encoder] (pool) {Attention + mean + max pooling};
        \node[block, below=of pool] (z) {Stage~1 flow embedding $\mathbf{z}^{(1)}_i$};
        \node[block, below=of z] (context) {Causal relation-specific context construction};
        \node[block, below=of context] (query) {Current-flow query over historical Stage~1 embeddings};
        \node[block, below=of query] (gate) {Gated context update, residual feed-forward network, and classifier};
        \node[block, below=of gate] (output) {Context-aware malicious-flow probability};

        \draw[arrow] (packet) -- (projection);
        \draw[arrow] (projection) -- (film);
        \draw[arrow] (flow) -- (film);
        \draw[arrow] (film) -- (encoder);
        \draw[arrow] (encoder) -- (pool);
        \draw[arrow] (pool) -- (z);
        \draw[arrow] (z) -- (context);
        \draw[arrow] (context) -- (query);
        \draw[arrow] (query) -- (gate);
        \draw[arrow] (gate) -- (output);
    \end{tikzpicture}
    \caption{Detailed architecture of the proposed two-stage detector. The
    statistical branch conditions Stage~1 packet tokens, while Stage~2 uses the
    target flow as the query and historical flows as keys and values.}
    \label{fig:method_architecture}
\end{figure}

\subsection{Stage~1: Intra-Flow Time-Aware Representation}
\label{subsec:method_stage1}

Stage~1 transforms a variable-length packet sequence into one flow embedding.
Its principal configuration consists of a record projection, time-aware
sinusoidal encoding, flow-conditioned token modulation, a pre-normalized
Transformer encoder, hybrid mask-aware pooling, and a two-class prediction
head.

\subsubsection{Packet-Record Projection}
\label{subsubsec:packet_projection}

For packet position $t$, the preprocessed feature vector
$\mathbf{x}_{i,t}\in\mathbb{R}^{d_p}$ is projected to the model dimension:

\begin{equation}
    \mathbf{e}_{i,t}
    =\operatorname{LN}
    \left(W_p\mathbf{x}_{i,t}+\mathbf{b}_p\right),
    \qquad
    \mathbf{e}_{i,t}\in\mathbb{R}^{d}.
    \label{eq:record_projection}
\end{equation}

The implementation permits an MLP projection, but the principal configuration
uses one linear projection followed by layer normalization. Categorical packet
fields have already been one-hot encoded, numerical fields have been transformed
and robustly scaled, and binary fields remain in $\{0,1\}$ as described in
Section~\ref{sec:dataset_preprocessing}.

\subsubsection{Cumulative Time-Aware Encoding}
\label{subsubsec:time_aware_encoding}

Conventional sinusoidal positional encoding represents packet order but cannot
distinguish two adjacent packets separated by microseconds from two adjacent
packets separated by seconds. Stage~1 therefore combines ordinal position with
cumulative elapsed time. The preprocessing pipeline stores

\begin{equation}
    u_{i,t}=\log\left(1+\Delta\tau_{i,t}\right),
    \label{eq:stored_log_iat}
\end{equation}

where $\Delta\tau_{i,t}$ is the nonnegative microsecond inter-arrival time from
the previous packet in the same bidirectional flow. Inside the model, the
interval is reconstructed and masked:

\begin{equation}
    \widehat{\Delta\tau}_{i,t}
    =m_{i,t}\max\left(0,\exp(u_{i,t})-1\right).
    \label{eq:reconstructed_iat}
\end{equation}

The cumulative elapsed time and its smoothed logarithm are

\begin{align}
    c_{i,t} &= \sum_{k=1}^{t}\widehat{\Delta\tau}_{i,k},\\
    p_{i,t} &= \log\left(1+c_{i,t}+\alpha\right),
    \qquad \alpha=10^{-7}.
    \label{eq:cumulative_time_feature}
\end{align}

Let the zero-based packet position be $q_t=t-1$. The combined coordinate is
$r_{i,t}=q_t+p_{i,t}$. For channel pair $j$, the time-aware encoding is

\begin{align}
    \operatorname{TE}_{i,t,2j}
    &=\sin\left(
        \frac{r_{i,t}}{10000^{2j/d}}
      \right),\\
    \operatorname{TE}_{i,t,2j+1}
    &=\cos\left(
        \frac{r_{i,t}}{10000^{2j/d}}
      \right).
    \label{eq:time_aware_sinusoid}
\end{align}

Padding positions receive a zero encoding. The encoded packet token is

\begin{equation}
    \widetilde{\mathbf{e}}_{i,t}
    =\operatorname{Dropout}
    \left(
        \mathbf{e}_{i,t}+m_{i,t}\operatorname{TE}_{i,t}
    \right).
    \label{eq:time_encoded_packet}
\end{equation}

This construction follows the time-aware Transformer principle used in GTID
\cite{Han2023GTID}, but it is applied here to structured, payload-content-free
packet metadata. Position-only and time-only modes are retained as ablations.
In the time-only mode, $r_{i,t}=p_{i,t}$; in the position-only mode,
$r_{i,t}=q_t$.

\subsubsection{Flow-Statistics Encoder}
\label{subsubsec:flow_encoder}

The separately preprocessed statistical flow vector $\mathbf{s}_i$ is mapped
to the packet-model dimension through a bottleneck encoder:

\begin{equation}
    \mathbf{f}_i
    =\operatorname{LN}_2
    \left(
        W_{f,2}\,
        \operatorname{Dropout}
        \left(
            \operatorname{GELU}
            \left(
                \operatorname{LN}_1(W_{f,1}\mathbf{s}_i+\mathbf{b}_{f,1})
            \right)
        \right)
        +\mathbf{b}_{f,2}
    \right).
    \label{eq:flow_feature_encoder}
\end{equation}

The principal bottleneck dimension is 32 and the output dimension is 128. The
bottleneck limits the capacity of the aggregate-feature branch and reduces the
risk that high-dimensional flow statistics overwhelm packet-sequence evidence.

\subsubsection{Bounded Flow-Conditioned Token Modulation}
\label{subsubsec:token_film}

The principal Stage~1 fusion mode is inspired by feature-wise linear modulation
(FiLM), in which conditioning information produces feature-specific affine
parameters \cite{Perez2018FiLM}. From the flow representation, a small network
generates

\begin{equation}
    [\boldsymbol{\gamma}_i;\boldsymbol{\beta}_i]
    =G_f(\mathbf{f}_i),
    \qquad
    \boldsymbol{\gamma}_i,\boldsymbol{\beta}_i\in\mathbb{R}^{d}.
    \label{eq:film_parameters}
\end{equation}

For each packet token, the raw flow-conditioned residual is

\begin{equation}
    \mathbf{r}^{\mathrm{raw}}_{i,t}
    =\eta
    \left(
        \boldsymbol{\gamma}_i\odot\widetilde{\mathbf{e}}_{i,t}
        +\boldsymbol{\beta}_i
    \right),
    \qquad
    \eta=\sigma(a),quad a_{0}=-5,
    \label{eq:film_raw_residual}
\end{equation}

where $\odot$ denotes element-wise multiplication. The initial scale is
$\sigma(-5)\approx0.0067$. In addition, the residual is constrained relative
to the norm of its packet token. With the configured maximum residual ratio
$\rho=0.13$,

\begin{align}
    \lambda_{i,t}
    &=\min\left(
        1,
        \frac{\rho\|\widetilde{\mathbf{e}}_{i,t}\|_2}
             {\|\mathbf{r}^{\mathrm{raw}}_{i,t}\|_2+\epsilon}
      \right),\\
    \mathbf{e}^{F}_{i,t}
    &=\widetilde{\mathbf{e}}_{i,t}
      +\lambda_{i,t}\mathbf{r}^{\mathrm{raw}}_{i,t}.
    \label{eq:bounded_token_film}
\end{align}

The final layer producing $\boldsymbol{\gamma}_i$ and
$\boldsymbol{\beta}_i$ is initialized to zero. The model therefore begins as
the packet-only branch and learns a bounded statistical correction. This is
particularly useful when aggregate statistics are predictive but potentially
noisy or capture-specific.

The terminology is important. Because the statistical vector is encoded in a
separate branch but modulates packet tokens before self-attention, the principal
configuration is most precisely described as \emph{separate-branch conditional
intermediate fusion}. It is not strict post-pooling late fusion. A post-pooling
gated fusion implementation is available as an alternative, but the reported
principal Scheme~C configuration uses the bounded token-FiLM mechanism in
Equation~\ref{eq:bounded_token_film}.

\subsubsection{Masked Transformer Encoder}
\label{subsubsec:stage1_transformer}

The conditioned sequence
$\mathbf{E}^{F}_i\in\mathbb{R}^{L\times d}$ is processed by a stack of
pre-normalized Transformer encoder layers. For layer $\ell$,

\begin{align}
    \mathbf{A}^{\ell}_i
    &=\operatorname{MHA}
      \left(
        \operatorname{LN}(\mathbf{H}^{\ell-1}_i),
        \operatorname{LN}(\mathbf{H}^{\ell-1}_i),
        \operatorname{LN}(\mathbf{H}^{\ell-1}_i);
        \mathbf{m}_i
      \right),\\
    \overline{\mathbf{H}}^{\ell}_i
    &=\mathbf{H}^{\ell-1}_i
      +\operatorname{Dropout}(\mathbf{A}^{\ell}_i),\\
    \mathbf{H}^{\ell}_i
    &=\overline{\mathbf{H}}^{\ell}_i
      +\operatorname{FFN}
       \left(\operatorname{LN}(\overline{\mathbf{H}}^{\ell}_i)\right),
    \label{eq:stage1_transformer_layer}
\end{align}

where $\mathbf{H}^{0}_i=\mathbf{E}^{F}_i$. Each attention head uses scaled
dot-product attention

\begin{equation}
    \operatorname{Attn}(Q,K,V)
    =\operatorname{softmax}
      \left(
        \frac{QK^{\top}}{\sqrt{d_h}}+M_i
      \right)V,
    \label{eq:scaled_dot_product_attention}
\end{equation}

where $M_i$ assigns negative infinity to padded key positions. The feed-forward
network uses GELU activation. The principal model has two encoder layers, four
attention heads, model dimension 128, feed-forward dimension 256, and dropout
0.15.

No causal mask is used inside Stage~1: all retained packets of the completed or
truncated flow can attend to one another. Consequently, Stage~1 is a
flow-level classifier, not a claim of packet-by-packet early detection. Padding
is excluded from attention keys and from the subsequent pooling operation.

\subsubsection{Hybrid Mask-Aware Pooling}
\label{subsubsec:hybrid_pooling}

The packet states are reduced using three complementary summaries. First, a
learned attention score is computed for each valid token:

\begin{align}
    a_{i,t}
    &=\mathbf{w}_{a,2}^{\top}
      \operatorname{GELU}
      \left(W_{a,1}\mathbf{h}_{i,t}+\mathbf{b}_{a,1}\right)+b_{a,2},\\
    \omega_{i,t}
    &=\frac{m_{i,t}\exp(a_{i,t})}
      {\sum_{k=1}^{L}m_{i,k}\exp(a_{i,k})},\\
    \mathbf{z}^{\mathrm{att}}_i
    &=\sum_{t=1}^{L}\omega_{i,t}\mathbf{h}_{i,t}.
    \label{eq:stage1_attention_pooling}
\end{align}

The masked mean and maximum are

\begin{align}
    \mathbf{z}^{\mathrm{mean}}_i
    &=\frac{\sum_{t=1}^{L}m_{i,t}\mathbf{h}_{i,t}}
            {\max(1,\sum_{t=1}^{L}m_{i,t})},\\
    \mathbf{z}^{\mathrm{max}}_i
    &=\max_{t:m_{i,t}=1}\mathbf{h}_{i,t}.
    \label{eq:mean_max_pooling}
\end{align}

The final Stage~1 embedding combines all three:

\begin{equation}
    \mathbf{z}^{(1)}_i
    =\operatorname{Dropout}
      \left(
        \operatorname{GELU}
        \left(
          \operatorname{LN}
          \left(
            W_o
            [\mathbf{z}^{\mathrm{att}}_i;
             \mathbf{z}^{\mathrm{mean}}_i;
             \mathbf{z}^{\mathrm{max}}_i]
            +\mathbf{b}_o
          \right)
        \right)
      \right).
    \label{eq:hybrid_pooling}
\end{equation}

Attention pooling can emphasize a small number of informative packets, mean
pooling represents overall behavior, and max pooling preserves strong local
activations. Their concatenation is projected back to dimension $d$ so Stage~2
receives one consistently sized vector per flow.

The Stage~1 class logits are

\begin{equation}
    \boldsymbol{\ell}^{(1)}_i
    =h_1(\mathbf{z}^{(1)}_i),
    \qquad
    \mathbf{p}^{(1)}_i
    =\operatorname{softmax}(\boldsymbol{\ell}^{(1)}_i).
    \label{eq:stage1_classifier}
\end{equation}

The exported embedding is taken before the final two-class layer. It therefore
retains information useful for contextual refinement rather than reducing each
flow to a single probability.

\subsubsection{Stage~1 Fusion Variants}
\label{subsubsec:stage1_variants}

The implementation supports the three controlled variants in
Table~\ref{tab:stage1_fusion_variants}. They share the same packet features,
sequence policy, split, Transformer capacity, training procedure, and classifier
unless an experiment explicitly states otherwise.

\begin{table}[t]
    \centering
    \caption{Stage~1 packet/flow integration variants.}
    \label{tab:stage1_fusion_variants}
    \small
    \begin{tabular}{@{}p{1.1cm}p{3.1cm}p{8.9cm}@{}}
        \toprule
        Scheme & Input structure & Integration mechanism \\
        \midrule
        A & Packet + repeated flow statistics & The transformed flow vector is
        concatenated to every packet before record projection. \\
        B & Packet only & The flow-statistics branch is disabled; this isolates
        packet order and time-aware encoding. \\
        C & Packet + separate flow branch & The flow vector is encoded once and
        conditionally modulates packet tokens through bounded token-FiLM. This
        is the principal proposed Stage~1 configuration. \\
        \bottomrule
    \end{tabular}
\end{table}

\subsection{Stage~2: Causal Inter-Flow Context Modelling}
\label{subsec:method_stage2}

Stage~2 asks whether a target flow should be interpreted differently when
recent related behavior is available. It operates only on Stage~1 embeddings
and context metadata. Raw packet fields, flow statistics, and labels of
historical flows are not loaded into the Stage~2 input.

\subsubsection{Embedding Interface and Chronological Order}
\label{subsubsec:embedding_interface}

Stage~1 exports, for each split,

\begin{equation}
    \left(
      \operatorname{flow\_id}_i,
      y_i,
      \mathbf{z}^{(1)}_i,
      t_i^{0},u_i,v_i,\rho_i
    \right),
    \label{eq:stage1_export_tuple}
\end{equation}

where $t_i^{0}$ is the flow start time, $u_i$ and $v_i$ are source and
destination identifiers, and $\rho_i$ is the split. Stage~2 validates unique
flow identifiers, concatenates the split files, and stably sorts by
$(t_i^{0},\operatorname{flow\_id}_i)$. The embedding rows are reordered with
the metadata so their alignment is preserved.

Training embeddings are exported with a deterministic, non-sampled loader.
This is essential because the weighted training sampler draws with replacement;
using it for export would duplicate some flows and omit others. The weighted
sampler is used only for optimization batches.

\subsubsection{Relation-Specific Causal Contexts}
\label{subsubsec:context_relations}

The context builder scans the sorted flows once from past to future. Before
adding target flow $i$ to any history, it retrieves eligible earlier indices.
The four evaluated relations are defined in
Table~\ref{tab:context_relations}.

\begin{table}[t]
    \centering
    \caption{Stage~2 context relations.}
    \label{tab:context_relations}
    \small
    \begin{tabular}{@{}p{3.0cm}p{4.0cm}p{7.1cm}@{}}
        \toprule
        Method & Eligibility condition for $j<i$ & Interpretation \\
        \midrule
        \methodfield{time_only} & Always eligible & Most recent flows globally,
        regardless of endpoint identity. \\
        \methodfield{source_host} & $u_j=u_i$ & Earlier flows initiated by the
        same source identifier. \\
        \methodfield{destination_host} & $v_j=v_i$ & Earlier flows directed to
        the same destination identifier. \\
        \methodfield{endpoint} &
        $\{u_j,v_j\}\cap\{u_i,v_i\}\neq\varnothing$ & Earlier flows sharing at
        least one endpoint in either role. \\
        \bottomrule
    \end{tabular}
\end{table}

For relation $r$, let

\begin{equation}
    \mathcal{A}^{(r)}_i
    =\left\{j<i:R_r(i,j)=1\right\}.
    \label{eq:method_eligible_context}
\end{equation}

After split-policy filtering and deduplication, the most recent $W-1$ indices
are retained:

\begin{equation}
    \mathcal{C}^{(r)}_i
    =\operatorname{Recent}_{W-1}
      \left(\mathcal{A}^{(r)}_i\right).
    \label{eq:recent_context}
\end{equation}

The current flow is then appended, producing
Equation~\ref{eq:method_stage2_input}. Histories are updated only after this
sequence has been built, which prevents self-inclusion and future leakage. The
principal context policy is \methodfield{online}: validation flows may use
earlier training flows, and test flows may use earlier training, validation,
and test flows, but never a future flow. This simulates persistent detector
state across chronological deployment. No historical label is exposed.

The term \methodfield{time_only} means chronological global recency. It does
not mean that continuous inter-flow time gaps are supplied to Stage~2. In the
current implementation, Stage~2 encodes relative context age, while explicit
microsecond timing is modelled inside Stage~1.

The collate function left-pads short sequences so the target remains at the last
valid position:

\begin{equation}
    [\mathrm{PAD},\ldots,\mathrm{PAD},
      \mathbf{z}^{(1)}_{j_1},\ldots,
      \mathbf{z}^{(1)}_{j_{K_i}},
      \mathbf{z}^{(1)}_i].
    \label{eq:left_padded_context}
\end{equation}

The accompanying Boolean mask makes the physical padding positions irrelevant
to positional encoding and attention.

\subsubsection{Context-Age Positional Encoding}
\label{subsubsec:context_age_encoding}

If a target has $K_i$ historical tokens, the target receives age 0, the most
recent historical flow receives age 1, and the oldest retained flow receives
age $K_i$. Let $a_{i,k}$ denote this age index. Stage~2 adds a conventional
sinusoidal encoding

\begin{align}
    \operatorname{PE}_{i,k,2j}
    &=\sin\left(
        \frac{a_{i,k}}{10000^{2j/d}}
      \right),\\
    \operatorname{PE}_{i,k,2j+1}
    &=\cos\left(
        \frac{a_{i,k}}{10000^{2j/d}}
      \right).
    \label{eq:stage2_age_encoding}
\end{align}

Unlike absolute tensor positions, age values are invariant to the amount of
left padding. This allows the same target and history order to receive the same
encoding in batches with different sequence lengths.

\subsubsection{Target-Query Gated Multi-Head Attention}
\label{subsubsec:target_query_attention}

The principal Stage~2 model is
\methodfield{target_query_gated}. Rather than allowing every historical token
to attend to every other historical token, it uses the current flow as one
query over the eligible history. Let $\mathbf{c}_i$ be the projected current
Stage~1 embedding and let
$\widetilde{\mathbf{H}}_i\in\mathbb{R}^{K_i\times d}$ contain age-encoded
historical embeddings. For attention head $h$,

\begin{align}
    \mathbf{q}^{h}_i
    &=W_Q^{h}\operatorname{LN}
      \left(\mathbf{c}_i+\operatorname{PE}(0)\right),\\
    \mathbf{K}^{h}_i
    &=\operatorname{LN}(\widetilde{\mathbf{H}}_i)W_K^{h},\\
    \mathbf{V}^{h}_i
    &=\operatorname{LN}(\widetilde{\mathbf{H}}_i)W_V^{h},\\
    \mathbf{h}^{h}_i
    &=\operatorname{softmax}
      \left(
        \frac{\mathbf{q}^{h}_i(\mathbf{K}^{h}_i)^{\top}}
             {\sqrt{d_h}}+M_i
      \right)
      \mathbf{V}^{h}_i.
    \label{eq:target_query_attention}
\end{align}

The head outputs are concatenated and projected to obtain the historical
summary $\mathbf{h}_i$. If no eligible historical flow exists,
$\mathbf{h}_i=\mathbf{0}$. This explicit fallback prevents a fabricated
padding token from contributing context.

The model compares the current representation and context summary through
difference and multiplicative interactions:

\begin{align}
    \mathbf{r}_i
    &=[\mathbf{c}_i;\mathbf{h}_i;
       \mathbf{c}_i-\mathbf{h}_i;
       \mathbf{c}_i\odot\mathbf{h}_i],\\
    \mathbf{v}_i
    &=[\mathbf{h}_i;
       \mathbf{c}_i-\mathbf{h}_i;
       \mathbf{c}_i\odot\mathbf{h}_i].
    \label{eq:target_context_interactions}
\end{align}

A dimension-wise gate and context update are then computed:

\begin{align}
    \mathbf{g}_i&=\sigma(G(\mathbf{r}_i)),\\
    \mathbf{d}_i&=U(\mathbf{v}_i),\\
    \mathbf{f}_i&=\operatorname{LN}
      \left(\mathbf{c}_i+\mathbf{g}_i\odot\mathbf{d}_i\right).
    \label{eq:gated_context_update}
\end{align}

When no history exists, both $\mathbf{g}_i$ and $\mathbf{d}_i$ are set to zero,
so the model reduces to the current-flow representation. Finally, a residual
feed-forward block produces

\begin{align}
    \mathbf{z}^{(2)}_i
    &=\operatorname{LN}
      \left(
        \mathbf{f}_i+\operatorname{FFN}(\mathbf{f}_i)
      \right),\\
    \boldsymbol{\ell}^{(2)}_i
    &=h_2(\mathbf{z}^{(2)}_i).
    \label{eq:stage2_output}
\end{align}

The target-query formulation preserves a strong current-flow path and lets
context act as a learned correction. Its attention-score matrix grows linearly
with $W$ because it contains one query, whereas full self-attention constructs
an $W\times W$ score matrix. This makes a window of 128 flows practical while
retaining multi-head retrieval from historical behavior.

\subsubsection{Stage~2 Control Models}
\label{subsubsec:stage2_controls}

To distinguish gains from context selection from gains caused merely by an
additional classifier, all control models consume the same Stage~1 embeddings,
context indices, left-padding convention, train/validation/test partitions,
and optimization pipeline.

\begin{itemize}[leftmargin=*]
    \item \textbf{No-context MLP} selects only the final valid target token and
    applies the same projection/classification style. It is the direct control
    for whether historical flows add useful information.

    \item \textbf{Vanilla Stage~2 Transformer} applies self-attention to the
    complete context sequence and classifies the last valid token. It tests
    whether all-to-all context interaction is preferable to target-query
    retrieval.

    \item \textbf{LSTM and GRU} use packed, right-compacted versions of the same
    context sequences and classify the last recurrent state
    \cite{Hochreiter1997LSTM,Cho2014GRU}.

    \item \textbf{CNN+LSTM} applies parallel one-dimensional convolutions with
    kernel sizes 3, 5, and 7, projects their concatenated outputs, and then uses
    an LSTM. This control combines local pattern extraction with recurrent
    sequence modelling \cite{Kim2014CNN,Hochreiter1997LSTM}.
\end{itemize}

These models are comparison baselines rather than components simultaneously
active in the proposed target-query model.

\subsection{Imbalance-Aware Sequential Training}
\label{subsec:method_training}

\subsubsection{Focal Objective and Controlled Sampling}
\label{subsubsec:focal_training}

Both stages use two-class softmax outputs and focal loss. For a sample with
label $y_i$ and predicted probability $p_{i,y_i}$,

\begin{equation}
    \mathcal{L}_{\mathrm{focal}}
    =-\frac{1}{B}
      \sum_{i=1}^{B}
      \alpha_{y_i}
      \left(1-p_{i,y_i}\right)^{\gamma}
      \log p_{i,y_i}.
    \label{eq:method_focal_loss}
\end{equation}

Focal loss reduces the influence of easily classified examples and was
originally proposed for severe foreground/background imbalance
\cite{Lin2017FocalLoss}. The principal configuration uses $\gamma=1$ and no
label smoothing. Stage~1 applies mild class weights $[1,1.25]$; Stage~2 uses
near-uniform weights $[1,1.01]$.

Training batches are additionally formed by a controlled weighted sampler. If
$n_0$ and $n_1$ are the training counts and the desired expected positive
fraction is $q$, each negative and positive example receives

\begin{equation}
    w_0=\frac{1-q}{n_0},
    \qquad
    w_1=\frac{q}{n_1}.
    \label{eq:controlled_sampler}
\end{equation}

Sampling is performed with replacement for exactly the original number of
training examples per epoch. The principal values are $q=0.09$ for Stage~1 and
$q=0.05$ for Stage~2. These values increase minority exposure relative to the
natural 3.945\% prevalence without forcing an artificial 50/50 training
distribution. Validation and test loaders never use weighted sampling.

\subsubsection{Optimization and Model Selection}
\label{subsubsec:optimization}

Both stages are optimized with AdamW, which decouples weight decay from the
adaptive gradient update \cite{Loshchilov2019AdamW}. Stage~1 uses learning rate
$10^{-4}$, weight decay $10^{-4}$, batch size 128, gradient-norm clipping at
1.0, and mixed-precision training when CUDA is available. A
ReduceLROnPlateau scheduler monitors validation PR-AUC, halves the learning rate
after five unimproved epochs, and does not reduce it below $10^{-6}$. Training
stops after 25 epochs without improvement, up to a maximum of 500 epochs.

Stage~2 uses learning rate $3\times10^{-5}$, weight decay
$3\times10^{-4}$, batch size 128, gradient-norm clipping at 0.5, and mixed
precision on CUDA. The first five epochs linearly warm the learning rate from
5\% to 100\% of its configured value. Cosine annealing then reduces it to 1\%
of the initial rate. Early stopping patience is 15 epochs with a maximum of 500
epochs.

The best checkpoint in each stage is selected by validation PR-AUC. This
threshold-independent criterion is appropriate for the imbalanced data and
avoids selecting a checkpoint from accuracy dominated by label~0
\cite{SaitoRehmsmeier2015}. The test partition is not used to choose the model.

\subsubsection{Sequential Stage Training and Embedding Export}
\label{subsubsec:sequential_training}

The complete training procedure is:

\begin{enumerate}[label=\textbf{Step \arabic*:},leftmargin=*]
    \item Fit the Stage~1 preprocessor only on the training partition and
    transform all partitions with the frozen preprocessing state.

    \item Train Stage~1 using the weighted training loader. Select the best
    checkpoint by validation PR-AUC and restore that checkpoint.

    \item Export exactly one Stage~1 embedding per flow using deterministic,
    non-sampled loaders. Store the flow identifier, label, split, start time,
    source identifier, and destination identifier with the embedding.

    \item Concatenate and chronologically sort the Stage~1 outputs. Construct
    causal context indices using only metadata and prior indices.

    \item Train Stage~2 while treating Stage~1 embeddings as fixed inputs.
    Select and restore the best Stage~2 checkpoint using validation PR-AUC.

    \item Select a decision threshold on validation predictions and apply that
    fixed threshold once to the test predictions.
\end{enumerate}

The current implementation does not back-propagate Stage~2 gradients into
Stage~1. Sequential training makes the source of an improvement auditable: a
Stage~2 gain is attributable to contextual modelling over a fixed representation,
not to an unreported change in the packet encoder.

\subsubsection{Validation-Selected Decision Threshold}
\label{subsubsec:decision_threshold}

For class-1 probability $p_i$, the prediction is

\begin{equation}
    \widehat{y}_i=\mathbb{I}[p_i\geq\delta^{*}].
    \label{eq:method_decision_rule}
\end{equation}

The principal implementation selects $\delta^{*}$ by maximizing class-1 F1 on
the validation partition. Stage~2 searches 199 thresholds between 0.01 and
0.99. The selected threshold is then frozen for test evaluation. Macro-F1 is
the primary reported aggregate F1 metric, while class-1 precision, recall, and
F1 are reported separately to make the threshold trade-off explicit. A
test-optimized threshold may be computed only as a clearly labelled diagnostic
oracle and must not be reported as deployable test performance. This separation
follows standard guidance against test-set adaptation in security machine
learning \cite{Arp2022DosDonts}.

\subsection{Principal Configuration}
\label{subsec:principal_configuration}

Table~\ref{tab:principal_method_config} records the architecture used for the
principal experiments. Values changed by an ablation, such as packet sequence
length, truncation strategy, context relation, or Stage~2 encoder type, must be
reported with that experiment rather than silently replacing this table.

\begin{table}[t]
    \centering
    \caption{Principal architecture and optimization configuration.}
    \label{tab:principal_method_config}
    \small
    \begin{tabular}{@{}p{2.0cm}p{4.4cm}p{7.6cm}@{}}
        \toprule
        Stage & Component & Principal value \\
        \midrule
        Stage~1 & Packet sequence & $L=128$, head selection, right padding \\
        Stage~1 & Packet encoder & $d=128$, 4 heads, 2 Transformer layers,
        FFN dimension 256, dropout 0.15 \\
        Stage~1 & Time encoding & Position + cumulative log-time,
        $\alpha=10^{-7}$ \\
        Stage~1 & Statistical branch & Bottleneck dimension 32, token-FiLM,
        initial scale logit $-5$, residual cap $\rho=0.13$ \\
        Stage~1 & Pooling & Learned attention + masked mean + masked max \\
        Stage~1 & Optimization & AdamW, $10^{-4}$ learning rate,
        $10^{-4}$ weight decay, batch 128, focal $\gamma=1$,
        sampled positive fraction 0.09 \\
        \midrule
        Stage~2 & Main encoder & Target-query gated multi-head attention \\
        Stage~2 & Context & Source-host in the principal/best configuration;
        time-only, destination-host, and endpoint are controlled comparisons \\
        Stage~2 & Context policy & $W=128$ including target, online causal
        history, deduplication, left padding \\
        Stage~2 & Representation & $d=128$, 4 heads, age positional encoding,
        FFN dimension 256, dropout 0.25 \\
        Stage~2 & Optimization & AdamW, $3\times10^{-5}$ learning rate,
        $3\times10^{-4}$ weight decay, batch 128, focal $\gamma=1$,
        sampled positive fraction 0.05 \\
        Both & Reproducibility & Random seed 130; validation PR-AUC checkpoint
        selection; validation-only threshold selection \\
        \bottomrule
    \end{tabular}
\end{table}

\subsection{Computational Properties}
\label{subsec:method_complexity}

For one Stage~1 flow, the self-attention score computation in each layer has
time and memory complexity $O(L^2d)$ and $O(L^2)$, respectively, while the
feed-forward block costs $O(Ld d_{\mathrm{ff}})$. The statistical encoder and
token-FiLM modulation are linear in $L$. Fixing $L=128$ therefore bounds both
memory and inference cost independently of the original packet count.

For Stage~2, context construction is performed once by maintaining bounded
global and host-indexed histories. The target-query attention map contains
$W-1$ scores per head rather than $W^2$ scores, giving attention-map time and
memory $O(Wd)$ and $O(W)$ per target, excluding the linear projections. The
bounded context window prevents high-activity hosts from producing unbounded
sequences.

The hierarchy also permits deployment choices that a monolithic model would
not. Stage~1 embeddings can be cached once per flow, and alternative Stage~2
relations or encoders can be evaluated without reprocessing packets. The cost
is that Stage~2 performance depends on the fixed Stage~1 representation and
that the current sequential implementation is not jointly optimized end to end.

\subsection{Methodological Summary}
\label{subsec:method_summary}

The proposed method first learns a packet-order- and time-aware flow
representation, conservatively conditions it on aggregate flow statistics, and
then refines the decision using causal histories of related flow embeddings.
The main Stage~1 contribution is the combination of cumulative time-aware
encoding, bounded flow-conditioned token modulation, and hybrid pooling. The
main Stage~2 contribution is a relation-controlled target-query attention
mechanism with a gated residual update and an exact no-history fallback. The
strict separation among predictors, context metadata, and labels, together
with chronological context construction and validation-only model selection,
keeps the method aligned with realistic temporal evaluation.
